\documentclass{article}
\usepackage{bur,calc,graphicx,multicol,shortvrb}
\usepackage[lowtilde]{url}
\usepackage[colorlinks=true,linkcolor=black,urlcolor=black,citecolor=BuRblue,breaklinks=true]{hyperref}
\usepackage[table]{xcolor}
\usepackage{tikz,dcolumn}
\urlstyle{rm}
\MakeShortVerb{|}
\Volume{1}
\Issue{1}
\Month{May}
\Year{08}

\newtheorem{definition}{Definition}
\newtheorem{lemma}{Lemma}


\begin{document}

\title{Please Place your Title Here:\newline
  Insert Line Breaks to Make it Look Nice}




\author{%
Author Name, \inst{Institute}, E-mail:
\and
Further Author Name, \inst{Further institute}, Further E-mail:
}

\begin{abstract}
  Place a short summary of the contents of your contribution here.
\end{abstract}

\keywords{Please enter some keywords here.}

\maketitle


  The text of your contribution starts here.

\section{Headings}


  Three levels of headings are supported. Please use |\section|, |\subsection| and |\subsubsection|.

\subsection{Heading level 2}

  This is an example of a |\subsection|.



\subsubsection{Heading level 3} (|\subsubsection|) will look like this.


  References may be cited using the |\cite| command as described, e.g., in \cite{mittelbach2005}. For more information about the layout of citations see subsection~\ref{citations}.


\section{Figures and Tables}

  Below you will find examples of tables. Please note and follow the special
  markup to achieve the desired BuR layout.  Some dummy text demonstrates the
  usage of figures and tables.

For tables use the |\tabularx| environment.  Additionally available columntypes are described in Table~\ref{tab:1}.


\begin{table}[htb!]
\small
\caption{Additional columntypes}\label{tab:1}
\noindent\begin{tabularx}{\columnwidth}{lL}
% use \normalsize and \textbf to highlight the table head
  \normalsize\textbf{Type}&
  \normalsize\textbf{Description}\\[0.5mm]\bhline
  \rowcolor[gray]{0.9}
    |C|
      &centered X column\\\hline
    |L|
      &raggedright X column\\\hline
    \rowcolor[gray]{0.9}
    |R|
      &raggedleft X column\\\hline
    |d|
      & allows vertical alignment of dots in numbers as shown in Table~\ref{tab:2}.
\end{tabularx}
\end{table}

\newcommand{\mc}[1]{
    \multicolumn{1}{c}{\normalsize\textbf{#1}}}

\begin{table}[htb!]
\rowcolors{1}{}{gray!0.9}
\small
\caption{Insert your caption here}\label{tab:2}
\noindent\begin{tabularx}{\columnwidth}{LCdX}
% use \normalsize and \textbf to highlight the table head
  \normalsize\textbf{magnim}&
  \normalsize\textbf{elit}&
  \mc{volortie}&% \\[0.5mm]\bhline at the last line of the table head
% is important to get the correct layout
  \normalsize\textbf{utat}\\[0.5mm]
% use rowcolor[gray]{0.9} to achieve alternating background colour of the table rows
  \rowcolor[gray]{0.9}
    tem &lor        &-48.345        &$\hphantom{0}12$\\\hline
    ullut   &sim        &-98.12        &$\hphantom{0}43$\\\hline
  \rowcolor[gray]{0.9}
    lortin  &corper     &-265.124       &$\hphantom{0}29$\\\hline
    del &dolorpe    &-50.0        &$122$\\\hline
  \rowcolor[gray]{0.9}
    tem &\multicolumn{3}{>{\columncolor[gray]{0.9}}c}{some text some text}\\\hline
        lortin  &corper     &-265.0       &$\hphantom{0}29$\\\hline
  \rowcolor[gray]{0.9}
    lortin  &corper     &-265.0       &$\hphantom{0}29$\\\hline
    del &dolorpe    &50.12        &$122$
\end{tabularx}
\legend{Add further comments to the table here\hfill\break
  Source: or a source line}
\end{table}



The background of table rows appears in white and gray by turns. To merge columns with gray background use

\small
\begin{verbatim}
\multicolumn{3}{%
    >{\columncolor[gray]{0.9}}columntype}
        {some text}
\end{verbatim}
\normalsize

as shown in row 5 in Table~\ref{tab:2}.

Wide tables are set in a |\table*| environment which is shown in Table~\ref{tab:3}.





% sample of two-column table
\begin{table*}[t!]
% you can use this definition inside tables to align numbers
\def~{\hphantom{0}}
\caption{Please insert your caption here}\label{tab:3} \small
% please use tabularx to gain tabulars that extend over column or text width, respectively
% note the X columns of tabularx
\noindent\begin{tabularx}{\textwidth}{XXXXXXXXX}
% use \normalsize and \textbf to highlight the table head
\normalsize\textbf{magnim}
  &\normalsize\textbf{elit}
  &\normalsize\textbf{volortie}
  &\normalsize\textbf{utat}
  &\normalsize\textbf{magnim}
  &\normalsize\textbf{elit}
  &\normalsize\textbf{volortie}
  &\normalsize\textbf{utat}
  &\normalsize\textbf{eppet}\\[0.5mm]\bhline
% use rowcolor[gray]{0.9} to achieve alternating background colour of the table rows
  \rowcolor[gray]{0.9}
    tem     &lor    &$~{-}48$    &$~12$  &tem     &lor     &$~{-}48$    &$~12$  &$~12$\\\hline
    ullut   &sim    &$~{-}98$    &$~43$  &ullut   &sim     &$~{-}98$    &$~43$  &$~43$\\\hline
  \rowcolor[gray]{0.9}
    lortin  &corper &  $-265$    &$~29$  &lortin  &corper  &$  -265$    &$~29$  &$~29$\\\hline
    del     &dolorpe&$~{-}50$    &$122$  &del     &dolorpe &$~{-}50$    &$122$  &$122$\\\hline
  \rowcolor[gray]{0.9}
    tem     &lor    &$~{-}48$    &$~12$  &tem     &lor     &$~{-}48$    &$~12$  &$~12$\\\hline
    ullut   &sim    &$~{-}98$    &$~43$  &ullut   &sim     &$~{-}98$    &$~43$  &$~43$\\\hline
  \rowcolor[gray]{0.9}
    lortin  &corper &  $-265$    &$~29$  &lortin  &corper  &$~{-}265$   &$~29$  &$~29$\\\hline
    del     &dolorpe&$~{-}50$    &$122$  &del     &dolorpe &$~{-}50$    &$122$  &$122$
\end{tabularx}
\legend{Add additional text corresponding to your table here.}
\end{table*}


Lorem ipsum dolor sit amet, consetetur sadipscing elitr, sed diam nonumy
  eirmod tempor invidunt ut labore et dolore magna aliquyam erat, sed diam
  voluptua.
\cite{loeffler1996a}
Lorem ipsum dolor sit amet, consetetur sadipscing elitr, sed diam nonumy
  eirmod tempor invidunt ut labore et dolore magna aliquyam erat, sed diam
  voluptua.

Lorem ipsum dolor sit amet, consetetur sadipscing elitr, sed diam nonumy
  eirmod tempor invidunt ut labore et dolore magna aliquyam erat, sed diam
  voluptua.

\begin{figure}[b!]
\caption{Place your caption here}
\begin{tikzpicture}[scale=2.5,domain=0.0:1,font=\rm]
    %\draw[very thin,color=gray] (-0.1,-1.1) grid (3.9,3.9);

\foreach \x/\xtext in {0.2, 0.4, 0.6, 0.8, 1.0}
    \draw[xscale=2,shift={(\x,0)}] (0pt,1pt) -- (0pt,-1pt) node[below] {$\xtext$};

  \foreach \y/\ytext in {-0.5/-5,0.5/5,1/10}
    \draw[yscale=1, shift={(0,\y)}] (1pt,0pt) -- (-1pt,0pt) node[left] {$\ytext$};


    \draw[->] (-0.2,0) -- (2.2,0) node[right] {$\tau$};
    \draw[->,yscale=1] (0,-0.6) -- (0,1.1) node[above] {$NPV / NPV^\tau$};
     \draw[color=black,xscale=2, yscale=1] plot[id=taxparadox3,smooth,samples=100] function{-300+(105)/1.1+(115)/(1.1)**2+145/(1.1)**3} node[right] {$NPV$};
    \draw[color=black,xscale=2, yscale=1] plot[id=taxparadox4,smooth,samples=100] function{-300+(105-x*(105-100))/(1+0.1*(1-x))+(115-x*(115-100))/(1+0.1*(1-x))**2+(145-x*(145-100))/(1+0.1*(1-x))**3} node[above] {$NPV^\tau$};
\end{tikzpicture}
\legend{Add additional text corresponding to your figure here.}
\end{figure}



Lorem ipsum dolor sit amet, consetetur sadipscing elitr, sed diam nonumy
  eirmod tempor invidunt ut labore et dolore magna aliquyam erat, sed diam
  voluptua. Lorem ipsum dolor sit amet, consetetur sadipscing elitr, sed diam nonumy
  eirmod tempor invidunt ut labore et dolore magna aliquyam erat, sed diam
  voluptua. Lorem ipsum dolor sit amet, consetetur sadipscing elitr, sed diam nonumy
  eirmod tempor invidunt ut labore et dolore magna aliquyam erat, sed diam
  voluptua.
\begin{figure}
  \heading{Video 1: Presentation of my First BuR Contribution}
\begin{tabularx}{\columnwidth}{|@{}C@{}|}
\hline
\includegraphics[width=\columnwidth-2pt]{Video-Sample.pdf} \\
\hline
\end{tabularx}
  \legend{Figure environments may be misused to place other objects with heading and/or caption}
\end{figure}


Lorem ipsum dolor sit amet, consetetur sadipscing elitr, sed diam nonumy
  eirmod tempor invidunt ut labore et dolore magna aliquyam erat, sed diam
  voluptua. Lorem ipsum dolor sit amet, consetetur sadipscing elitr, sed diam nonumy
  eirmod tempor invidunt ut labore et dolore magna aliquyam erat, sed diam
  voluptua. Lorem ipsum dolor sit amet, consetetur sadipscing elitr, sed diam nonumy
  eirmod tempor invidunt ut labore et dolore magna aliquyam erat, sed diam
  voluptua.

\begin{table}[b]
\small
\caption{For some tables alternating line colours are not appropriate}
% please use tabularx to gain tabulars that extend over column or text width, respectively
% note the X columns of tabularx
\noindent\begin{tabularx}{\columnwidth}{XcXXX}
% use \normalsize and \textbf to highlight the table head
  \normalsize\textbf{magnim}&&
  \normalsize\textbf{elit}&
  \normalsize\textbf{volortie}&% \\[0.5mm]\bhline at the last line of the table head
% is important to get the correct layout
  \normalsize\textbf{utat}\\[0.5mm]\bhline
    tem     &&lor        &$\hphantom{0}{-}48$        &$\hphantom{0}12$\\\hline
            &&sim        &$\hphantom{0}{-}98$        &$\hphantom{0}43$\\
    lortin  &$\smash{\left\{\vrule width0pt height6.5mm depth3.5mm\right.}$
            &corper     &$-265$       &$\hphantom{0}29$\\
            &&dolorpe    &$\hphantom{0}{-}50$        &$122$\\\hline
    tem     &&lor        &$\hphantom{0}{-}48$        &$\hphantom{0}12$\\\hline
    ullut   &&sim        &$\hphantom{0}{-}98$        &$\hphantom{0}43$\\
    lortin  &\smash{\raise2mm\hbox{$\left\{\vrule width0pt height4mm depth2mm\right.$}}
            &corper     &$-265$       &$\hphantom{0}29$\\\hline
    del     &&dolorpe    &$\hphantom{0}{-}50$        &$122$
\end{tabularx}
\legend{Add further comments to the table here\hfill\break
  Source: or a source line}
\end{table}

Lorem ipsum dolor sit amet, consetetur sadipscing elitr, sed diam nonumy
  eirmod tempor invidunt ut labore et dolore magna aliquyam erat, sed diam
  voluptua. Lorem ipsum dolor sit amet, consetetur sadipscing elitr, sed diam nonumy
  eirmod tempor invidunt ut labore et dolore magna aliquyam erat, sed diam
  voluptua.

\section{Math}

Business Research appears in twocolumn mode. This fact might sometimes lead to a challenging task when typesetting formulas. To provide a little more space for your formulas you might adjust the horizontal space between the formula label and the formula itself. Use in your preamble e.g.

\begin{equation}
E = m \times c^2
\end{equation}



\bigskip
\verb+\setlength{\mathindent}{3pt}+
\bigskip
\setlength{\mathindent}{3pt}

for that adjustment.

\begin{equation}
E = m \times c^2
\end{equation}






\section{Theorems}

Use the \texttt{theorem} environment for definitions, lemmas, hypothesis etc. To create individual desired theorems e.g. ``Definition'' use

\bigskip
\verb+\newtheorem{definition}{Definition}+
\bigskip

in your preamble. The theoremstyle \texttt{bur} then leads to the following result


\begin{definition}[Lorem ipsum dolor sit amet]
Lorem ipsum dolor sit amet, consetetur sadipscing elitr, sed diam nonumy
  eirmod tempor invidunt ut labore et dolore magna aliquyam erat, sed diam
  voluptua.
\end{definition}

\bigskip
\begin{verbatim}
\begin{definition}[Lorem ipsum ...]
Lorem ipsum dolor sit ...
\end{definition}
\end{verbatim}
\bigskip

Theorems are labeled and have an optional item in round brackets. The content of the theorem appears in italics.


\section{Citation style and bibliography style (with BiB\TeX)}

The following specifications have to be used for citations or the bibliography, respectively. In case of the citation layout, the examples are followed by the BiB\TeX\xspace syntax for BiB\TeX\xspace users. The bibliography entries are described in appendix~\ref{bib}.



The \texttt{bur.bst} file as the style file for the biblio\-graphy works fine. No adjustments of the \texttt{.bbl}-file are necessary. If you don't use BiB\TeX\xspace please be sure to provide your contribution with the specified layout. Note that because of the usage of the \texttt{natbib}-package the commands \verb+\cite+ and \verb+\citet+ lead to the same result.


\subsection{Citation style}\label{citations}

References have to be cited by the author-year system. All hints and comments are embedded in the continuous text. Please do not use footnotes for citations. The following examples illustrate the desired style for citations:

\begin{enumerate}
\item If citations are enclosed in round brackets, the corresponding year is not enclosed by round brackets additionally; page reference is located after the year followed by a colon:

Lorem ipsum dolor sit amet, consetetur sadipscing elitr, sed diam (\citealp[349]{schanz2011}).

  Syntax: \verb+(\citealp[349]{schanz2011})+


\item In case of more than one source, the sources are separated by a semicolon:

Lorem ipsum dolor sit amet, consetetur sadipscing elitr, sed diam (\citealp{bottazzi1998}; \citealp{fehr2003}).

Syntax:

\verb+(\citealp{bottazzi1998};+

\verb+\citealp{fehr2003})+




\item A little more work is caused by more than one source plus page references:

Lorem ipsum dolor sit amet, consetetur sadipscing elitr, sed diam (\citealp*[201-203]{bottazzi1998}; \citealp*[309 et seq.]{fehr2003}).

Syntax:

\verb+(\citealp*[201-203]{bottazzi1998};+
\verb+\citealp*[309 et seq.]{fehr2003}) +



\item If the source is embedded in the continuous text and not displayed in round brackets at the end of sentence use the following style

Lorem ipsum \citet{abdulkader2006} dolor sit amet, consetetur sadipscing elitr, sed diam.

Syntax: \verb+\citet{abdulkader2006}+

or

Lorem ipsum \citet[1701]{abdulkader2006} dolor sit amet, consetetur sadipscing elitr, sed diam.

Syntax: \verb+\citet[1701]{abdulkader2006}+

when page references are applied.

\item References with the same author but different years are cited as follows:

Lorem ipsum \citet{albers1996,albers2000} dolor sit amet, consetetur sadipscing elitr, sed diam.

Syntax: \verb+\citet{albers1996,albers2000}+

Lorem ipsum \citeauthor{albers1996} \citeblue{(}\citeyear{albers1996}, \citeyear{albers2000}, and \citeyear{albers2009}\citeblue{)} dolor sit amet, consetetur sadipscing elitr, sed diam.

Syntax:

\begin{verbatim}
\citeauthor{albers1996}
\citeblue{(}\citeyear{albers1996},
\citeyear{albers2000}, and
\citeyear{albers2009}\citeblue{)}
\end{verbatim}


\end{enumerate}


\subsection{Bibliography style}

The reference list should ideally be sorted in alphabetical order. References are provided without any indentation. The following description includes the most common types of references. Journal names are not abbreviated.

All words of a title have a leading capital letter except for articles (a, an, the), conjunctions and prepositions of up to four letters
(e.g. on, of, at, to, by, and, or, but, from, with).

The references are sorted by the last names of the authors whereas e.g. ``von'' is not considered as a part of the last name, ''von Krogh'' comes before ``Sureth''. The articles as for example ``van'' and ``della'' are part of the last name. That means that ``Sureth'' comes before ``Van Gogh''. To get the right sorting order please use \verb+author={Van Gogh, Vincent}+ or \verb+author={Vincent {Van Gogh}+ in your \verb+.bib+-Database.



\begin{enumerate}

\item \textbf{Monographs}

%(first edition is not mentioned, there is a comma followed by ``and'' before the last author if there are more than two authors)

\cite{kruschwitz2007}

Kruschwitz, Lutz (2007): \textit{Investitions\-rechnung}, 11th ed., Oldenbourg Verlag: M�nchen.


\cite{schanz2011}

Schanz, Deborah and Sebastian Schanz (2011): \textit{Business Taxation and Financial Decisions}, Springer: Heidelberg.


\cite{mittelbach2004}

Mittelbach, Frank, Michel Goossens, Johannes Braams, David Carlisle, and Chris Rowley (2004): \textit{The LaTeX Companion. Addison-Wesley Series on Tools and Techniques for Computer T}, 2nd ed., Addison-Wesley Professional: M�nchen.


\cite{cowley2008}

Cowley, Nick, Carlos Guti�rrez, Juhani Kesti and Mei-June Soo (eds.) (2008): \textit{Global Corporate Tax Handbook 2008}, International Bureau of Fiscal Documentation: Amsterdam.








\item \textbf{Journals}

\cite{loeffler1996b}

Andreas L�ffler (1996b): Variance Aversion Implies $\mu$-$\sigma^2$-Criterion, \textit{Journal of Economic Theory}, 69 (3): 532-539.


\cite{niemann2005}

Niemann, Rainer and Caren Sureth (2005): Capital Budgeting with Taxes under Uncertainty and Irreversibility, \textit{Journal of Economics and Statistics}, 225: 77-95.





\cite{krogh1996}

Krogh, Georg von and Johan Roos (1996): A Tale of the Unfinished, \textit{Strategie Management Journal}, 19 (9): 729-737.



\cite{arnold2011}

Arnold, Markus C. and Robert M. Gillenkirch (2001): Centralization Versus Delgation in an Experimental Capital Budgeting Setting, \textit{Business Research}, 4 (1): 10-30, \url{http://www.business-research.org/2011/1/accounting/2975/arnold-gillenkirch-centralization.pdf} (Access date: 2011-11-08).




\item \textbf{Incollections} (notice, name and surname of the first editor are not reversed)

\cite{fehr2003}

Fehr, Hans and Wolfgang Wiegard (2003): ACE for Germany? Fighting for a Better Tax System, in: Michael Ahlheim, Heinz-Dieter Wenzel, and Wolfgang Wiegard (eds.): \textit{Steuerpolitik \endash\xspace Von der Theorie zur Praxis: Festschrift f�r Manfred Rose}, Springer Verlag: Berlin et al., 297-324.

\cite{offermanns2006}

Offermanns, Ren� (2006): Belgium, in: Juhani Kesti (ed.), \textit{European Tax Handbook 2006}, International Bureau of Fiscal Documentation: Amsterdam, 63-82.



\item \textbf{Inproceedings}

\cite{matta2005}

Matta, Andrea and Robert Chefson (2005): Formal Properties of Closed Flow Lines with Limited Buffer Capacities and Random Processing Times, in: J. Manuel Felix-Teixera and A. E. Carvalho Brito (eds.): \textit{The 2005 European Simulation and Modelling Conference}, Porto, 190-198.






\item \textbf{Unpublished} (e.g. working paper, unpublished master thesis or internet resources)

\cite{hundsdoerfer2011}

Hundsdoerfer, Jochen, Christian Sielaff, Kay Blaufus,
Dirk Kiesewetter, and Joachim Weimann (2011): The Influence of Tax Labeling and Tax Earmarking on the Willingness to Contribute \endash\xspace A Conjoint Analysis, Working Paper, arqus, \url{http://www.arqus.info/mobile/paper/arqus_121.pdf} (Access date: 2011-09-28).


\cite{brown2000}

Brown, Lawrence D. (2000): I/B/E/S Research Bibliography, 6th ed., http:// www2.gsu.edu/\url{~}wwwaceFaculty/lbrown/Bibli ography.pdf (Access date: 2011-02-15).



\cite{overesch2006}

Overesch, Michael and Georg Wamser (2006): German Inbound Investment, Corporate Tax Planning, and Thin Capitalization Rules \endash\xspace a Difference-in-Difference Approach, Working Paper, Zentrum f�r Europ�ische Wirtschaftsforschung, \url{http://papers.ssrn.com/sol3/papers.cfm?abstract_id=945327} (Access date: 2011-11-18).


\cite{mintz2005}

Mintz, Jack and Alfons J. Weichenrieder (2005): Taxation and the Financial Structure of German Outbound FDI, Working Paper, CESifo, \url{http://papers.ssrn.com/sol3/papers.cfm?abstract_id=875306##} (Access date: 2011-11-18).


\cite{arnold2009}

Arnold, Markus C. and Dominik Schreiber (2009): Audits, Reputation,
and Repeated Interaction in a Capital Budgeting Setting, Working Paper, SSRN, \url{http://papers.ssrn.com/sol3/papers.cfm?abstract_id=1935469} (Access date: 2011-11-18).





\end{enumerate}

\section{Technical hints for the bibliography}

If you apply BiB\TeX\xspace the reference section is produced by

\begin{verbatim}
\bibliographystyle{bur}
\bibliography{bur}
\end{verbatim}

To run BiB\TeX\xspace properly, you have to provide a database that includes the cited references. Here the name of the bib-file has to be \texttt{bur.bib}.

When not using BiB\TeX\xspace use the following syntax for the reference section

\begin{verbatim}
\begin{thebibliography}{}
\item{Reference 1}
\item{Reference 2}
\end{thebibliography}
\end{verbatim}

In this case no separate database is needed. However, be sure that just the cited literature is included in the reference section.

%\cite{loeffler1996b}

%\citet*{bottazzi1998}

%\citet*{abdulkader2006}

%\citet*{beatty1986}

%\citet*{belsley1980}

%\citet*{botosan1997}

%\citet*{bradshaw2003}

%\citet*{brennan1970}

%\citet*{brown2000}

%\citet*{helber1999}

%\citet*{hopp2000}

%\citet*{matta2005}

%\citet*{smith1999}


%\citet*{fehr2003}

%\cite{mittelbach2005}


% Use the bibliography style for BuR
\bibliographystyle{bur}
% Use the bur.bib database to extract the references
\bibliography{bur}


\section*{Biographies}

\textbf{Alexandra Ma\ss baum} is doctoral student and research assistant at the Chair of
Business Administration, esp. Business Taxation at the University of Paderborn. She received her diploma degree in Business Administration with specialization in accounting, taxation, auditing, and management accounting, as well as corporate and tax law from the University of Osnabr�ck.


\medskip
\textbf{Caren Sureth} is Professor and Chair of Business Administration, esp. Business Taxation since 2004 and also Head of the Center for Tax Research at the Unversity of Paderborn. She received her doctorate from the University of Bielefeld, where she also reached her habilitation. She studied Business Administration, English, French, and Chinese and received her Diploma Degree at the University of Passau. Her main fields of research are the impact of taxation on entrepreneurial decisions, tax planning and taxation under uncertainty.




\onecolumn

\appendix

\section{BiB\TeX\xspace entries}\label{bib}

\begin{enumerate}
\item \textbf{Monographs}

\begin{verbatim}
@BOOK{kruschwitz2007,
  author    = {Lutz Kruschwitz},
  title     = {Investitionsrechnung},
  publisher = {Oldenbourg Verlag},
  address   = {M�nchen},
  edition   = {11},
  year      = {2007}
}
\end{verbatim}

\begin{verbatim}
@book{brown2000,
    author  = {Brown, Lawrence D.},
    title   = {I/B/E/S Research Bibliography},
    edition = {6th},
    year    = {2000},
    note    = {\url{http://www2.gsu.edu/~wwwaceFaculty/lbrown/Bibliography.pdf}
              (Access date: 2011-02-15)}
}
\end{verbatim}

\begin{verbatim}
@book{schanz2011,
    author    = {Schanz, Deborah and Sebastian Schanz},	
    title     = {Business Taxation and Financial Decisions},
    publisher = {Springer Verlag},
    address   = {Heidelberg},
    edition   = {},
    year      = {2011}
}
\end{verbatim}

\begin{verbatim}
@book{mittelbach2004,
    author    = {Mittelbach, Frank and Michel Goossens and Johannes Braams
                 and David Carlisle and Chris Rowley},
    title     = {The LaTeX Companion. Addison-Wesley Series on Tools
                 and Techniques for Computer T},
    publisher = {Addison-Wesley Professional},
    address   = {M�nchen},
    edition   = {2},
    year      = {2004}
}
\end{verbatim}

\begin{verbatim}
@book{cowley2008,
    editor={Cowley, Nick and Carlos Guti�rrez and Juhani Kesti and Mei-June Soo},	
    title={Global Corporate Tax Handbook 2008},
    publisher={International Bureau of Fiscal Documentation},
    address = {Amsterdam},
    edition={},
    year={2008}
}
\end{verbatim}

\item \textbf{Journals}


\begin{verbatim}
@article{loeffler1996b,
    author    = {L�ffler, Andreas},	
    title     = {Variance Aversion Implies $\mu$-$\sigma^2$-Criterion},
    journal   = {Journal of Economic Theory},
    volume    = {69},
    number    = {3},
    year      = {1996},
    pages     = {532-539}
}
\end{verbatim}

\begin{verbatim}
@article{niemann2005,
    author={Niemann, Rainer and Caren Sureth},
    title={Capital Budgeting with Taxes under Uncertainty and Irreversibility},
    journal={Journal of Economics and Statistics},
    volume={225},
    number = {},
    year={2005},
    pages={77-95}
}
\end{verbatim}


\begin{verbatim}
@article{krogh1996,
    author  = {Krogh, Georg von and Johan Roos},
    title   = {A Tale of the Unfinished},
    journal = {Strategie Management Journal},
    volume  = {19},
    number  = {9},
    year    = {1996},
    pages   = {729-737}
}
\end{verbatim}

\begin{verbatim}
@Article{arnold2011,
    author  = {Arnold, Markus C. and Robert M. Gillenkirch},
    title   = {Centralization Versus Delgation in an Experimental Capital Budgeting Setting},
    pages   = {10-30},
    year    = {2011},
    journal = {Business Research},
    volume  = {4},
    number  = {1},
    note    = {\url{http://www.business-research.org/2011/1/accounting/2975/
              arnold-gillenkirch-centralization.pdf} (Access date: 2011-11-08)}
}

\end{verbatim}




\item \textbf{Incollections}

\begin{verbatim}
@incollection{fehr2003,
    author     = {Hans Fehr and Wolfgang Wiegard},
    title      = {ACE for Germany? Fighting for a Better Tax System},
    year       = {2003},
    pages      = {297-324},
    editor     = {Michael Ahlheim and Heinz-Dieter Wenzel and Wolfgang Wiegard},
    booktitle  = {Steuerpolitik \endash\xspace Von der Theorie zur Praxis:
                  Festschrift f�r Manfred Rose},
    publisher  = {Springer},
    address    = {Berlin et al.}
}
\end{verbatim}

\begin{verbatim}
@incollection{offermanns2006,
    author     = {Ren� Offermanns},
    title      = {Belgium},
    year       = {2006},
    pages      = {63-82},
    editor     = {Juhani Kesti},
    booktitle  = {European Tax Handbook 2006},
    publisher  = {International Bureau of Fiscal Documentation},
    address    = {Amsterdam}
}
\end{verbatim}

\item \textbf{Inproceedings}

\begin{verbatim}
@inproceedings{matta2005,
    author     = {Matta, Andrea and Robert Chefson},
    title      = {Formal Properties of Closed Flow Lines with Limited Buffer
                  Capacities and Random Processing Times},
    year       = {2005},
    pages      = {190-198},
    editor     = {Felix-Teixera, J. Manuel and A. E. Carvalho Brito},
    booktitle  = {The 2005 European Simulation and Modelling Conference},
    address    = {Porto}
}
\end{verbatim}


\item \textbf{Unpublished}

\begin{verbatim}
@unpublished{hundsdoerfer2011,
 author = {Hundsdoerfer, Jochen and Christian Sielaff and Kay Blaufus
           and Dirk Kiesewetter and Joachim Weimann},
 title  = {The Influence of Tax Labeling and Tax Earmarking on the
           Willingness to Contribute \endash\xspace A Conjoint Analysis},
 note   = {Working Paper, arqus, \url{http://www.arqus.info/mobile/
           paper/arqus_121.pdf} (Access date: 2011-09-28)},
 year   = {2011}
}
\end{verbatim}

\begin{verbatim}
@unpublished{overesch2006,
 author    = {Overesch, Michael and Georg Wamser},
 title     = {German Inbound Investment, Corporate Tax Planning, and Thin
              Capitalization Rules \endash\xspace a Difference-in-Difference
              Approach},
 note      = {Working Paper, Zentrum f�r Europ�ische Wirtschaftsforschung,
              \url{http://papers.ssrn.com/sol3/papers.cfm?abstract_id=945327}
              (Access date: 2011-11-18)},
 year      = {2011}
}
\end{verbatim}

\begin{verbatim}
@unpublished{mintz2005,
 author    = {Mintz, Jack and Alfons J. Weichenrieder},
 title     = {Taxation and the Financial Structure of German Outbound FDI},
 note      = {Working Paper, CESifo, \url{http://papers.ssrn.com/sol3/
              papers.cfm?abstract_id=875306##} (Access date: 2011-11-18)},
 year      = {2011}
}
\end{verbatim}


\begin{verbatim}
@unpublished{arnold2005,
 author    = {Arnold, Markus C. and Dominik Schreiber},
 title     = {Audits, Reputation, and Repeated Interaction in a Capital
              Budgeting Setting},
 note      = {Working Paper, SSRN, \url{http://papers.ssrn.com/sol3/
              papers.cfm?abstract_id=1935469} (Access date: 2011-11-18)},
 year      = {2005}
}
\end{verbatim}



\end{enumerate}


\end{document}
